%; whizzy paragraph -pdf xpdf -latex ./whizzypdfptex.sh
%; whizzy-paragraph "^\\\\begin{frame}"
% latex beamer presentation.
% platex, latex-beamer でコンパイルすることを想定。 

%     Tokyo Debian Meeting resources
%     Copyright (C) 2009 Junichi Uekawa
%     Copyright (C) 2009 Nobuhiro Iwamatsu

%     This program is free software; you can redistribute it and/or modify
%     it under the terms of the GNU General Public License as published by
%     the Free Software Foundation; either version 2 of the License, or
%     (at your option) any later version.

%     This program is distributed in the hope that it will be useful,
%     but WITHOUT ANY WARRANTY; without even the implied warreanty of
%     MERCHANTABILITY or FITNESS FOR A PARTICULAR PURPOSE.  See the
%     GNU General Public License for more details.

%     You should have received a copy of the GNU General Public License
%     along with this program; if not, write to the Free Software
%     Foundation, Inc., 51 Franklin St, Fifth Floor, Boston, MA  02110-1301 USA

\documentclass[cjk,dvipdfmx,12pt]{beamer}
\usetheme{Tokyo}
\usepackage{monthlypresentation}

%  preview (shell-command (concat "evince " (replace-regexp-in-string
%  "tex$" "pdf"(buffer-file-name)) "&")) 
%  presentation (shell-command (concat "xpdf -fullscreen " (replace-regexp-in-string "tex$" "pdf"(buffer-file-name)) "&"))
%  presentation (shell-command (concat "evince " (replace-regexp-in-string "tex$" "pdf"(buffer-file-name)) "&"))

%http://www.naney.org/diki/dk/hyperref.html
%日本語EUC系環境の時
\AtBeginDvi{\special{pdf:tounicode EUC-UCS2}}
%シフトJIS系環境の時
%\AtBeginDvi{\special{pdf:tounicode 90ms-RKSJ-UCS2}}

\title{caff と mail-transport-agent}
\subtitle{第156回 2017年10月度 東京エリアDebian勉強会}
\author{yy\_y\_ja\_jp}
\date{2017年10月21日}
\logo{\includegraphics[width=8cm]{image200607/openlogo-light.eps}}

\begin{document}

\frame{\titlepage{}}

\section{Agenda}

\begin{frame}{Agenda}
\begin{enumerate}
 \item caff
 \item Mail Transport Agent
 \item mail-transport-agent
 \item caff と mail-transport-agent
\end{enumerate}
\end{frame}

\section{caff}
\begin{frame}{caff}
\begin{itemize}
 \item caff(1) -- CA - Fire \& Forget

 キーサインパーティなどで便利なツール

 \item 
 signing-party パッケージでインストールできる

 \item 
 キーサインパーティに参加した後で、このコマンドを実行すれば必要な作業をやってくれる

 \pause

 \begin{enumerate}
  \item 参加した人の公開鍵を取得
  \item その公開鍵に自分の秘密鍵で署名
  \item 署名済みの公開鍵をその人にメールで送信
 \end{enumerate}

\end{itemize}
  
\end{frame}

\section{Mail Transport Agent}
\begin{frame}{Mail Transport Agent (1/3)}
 メール転送エージェント MTA

 受け付けたメッセージを他のサーバなどに転送する
 \pause
 \begin{itemize}
  \item 受け付けかた（インターフェース）
	\begin{itemize}
	 \item SMTP（TCP 25番ポートなど）
	 \item sendmail互換（/usr/sbin/sendmailコマンドを実行して標準入力に
	       渡すなど）
	 \item など
	\end{itemize}
 \pause
  \item 転送のしかた
	\begin{itemize}
	 \item 自力でDNSを引いて送付先サーバに転送
	 \item 他のサーバに渡して転送してもらう（リレー）
	\end{itemize}
 \end{itemize}
 
\end{frame}

\begin{frame}{Mail Transport Agent (2/3)}
 受け付けかた（インターフェース）
	\begin{itemize}
	 \item SMTP（TCP 25番ポートなど）
	 \item sendmail互換（/usr/sbin/sendmailコマンドを実行して標準入力に
	       渡すなど）
	 \item など
	\end{itemize}

 メッセージをマシン上にある（ローカル）MTAに
	       渡すだけなら、わざわざメールサーバ（SMTPサーバ）を立ち上
 げてそれ経由で渡す必要はない\\
 -- /usr/sbin/sendmailコマンドを起動して、
 そのコマンド自体が転送してくれれば充分
\end{frame}

\begin{frame}{Mail Transport Agent (3/3)}
転送のしかた
	\begin{itemize}
	 \item 自力でDNSを引いて送付先サーバに転送
	 \item 他のサーバに渡して転送してもらう（リレー）
	\end{itemize}

 \pause
 自力で転送してしまうと
 \pause
 \begin{itemize}
  \item そのメールサーバをインターネット公開にしていない場合は自分宛てのメー
	ルが受信できない -- エラーメールも受け取れない
  \pause
  \item 自分のメールアドレスのドメインがそのサーバ管理でない場合はなりす
	まし（フィッシング）になってしまう

	\begin{quote}
This message may not have been sent by: xxx@example.com
	\end{quote}
\begin{quote}
Why is this message in Spam? It's similar to messages that were detected by our spam filters.
\end{quote}
  \pause
  \item OP25B の影響で送信できないこともある
 \end{itemize}

\end{frame}

\section{mail-transport-agent}
\begin{frame}[containsverbatim]{mail-transport-agent (1/4)}

\begin{itemize}
 \item Debianポリシーで定められた仮想パッケージのうちの1つ

 \scriptsize
 \verb|/usr/share/doc/debian-policy/virtual-package-names-list.txt.gz|
 \footnotesize
 \url{https://www.debian.org/doc/packaging-manuals/virtual-package-names-list.txt}
 \normalsize
 で見られる

\begin{quote}
mail-transport-agent{ }{ }{ }{ }a mail transport agent (e.g. Smail, Sendmail, \&c)
\end{quote}

 \item 
 mail-transport-agent を提供しているパッケージは
 \verb|apt-cache showpkg mail-transport-agent|
 や
 \footnotesize
 \url{https://packages.debian.org/stable/mail-transport-agent}
 \normalsize
 で見られる
\end{itemize}
\end{frame}

\begin{frame}[containsverbatim]{mail-transport-agent (2/4)}
 \begin{itemize}
  \item Debianのデフォルトは exim4-daemon-light\\
	-- 仮想パッケージ default-mta も提供(Provides)
 \item どのパッケージも Conflicts: mail-transport-agent, Replaces:
       mail-transport-agent, Provides: mail-transport-agent と書いてある\\
       ⇒ 2つ以上インストールできない
 \end{itemize}

 \scriptsize
\begin{verbatim}
$ apt-cache show exim4-daemon-light
Package: exim4-daemon-light
Source: exim4
Version: 4.89-2+deb9u1
Installed-Size: 1167
Maintainer: Exim4 Maintainers <pkg-exim4-maintainers@lists.alioth.debian.org>
Architecture: amd64
Replaces: exim4-base (<= 4.61-1), mail-transport-agent
Provides: default-mta, exim4-localscanapi-2.0, mail-transport-agent
Depends: exim4-base (>= 4.89), libc6 (>= 2.16), libdb5.3, libgnutls30 (>= 3.5.6), libpcre3, debconf (>= 0.5) | debconf-2.0
Conflicts: mail-transport-agent
\end{verbatim}
 \normalsize
\end{frame}

\begin{frame}[containsverbatim]{mail-transport-agent (3/4)}
 \begin{itemize}
 \item MTAのパッケージが2つに分かれていることもある --
       mail-transport-agentを提供するパッケージがsendmail互換インターフェー
       ス(/usr/sbin/sendmail)などを提供
 \end{itemize}

 \scriptsize
\begin{verbatim}
$ apt-cache show msmtp-mta
Package: msmtp-mta
Source: msmtp
Version: 1.6.6-1
（略）
Replaces: mail-transport-agent
Provides: mail-transport-agent
Depends: msmtp
Conflicts: mail-transport-agent
Description-en: light SMTP client with support for server profiles - the regular MTA
（以下略）
\end{verbatim}
\begin{verbatim}
$ apt-cache show msmtp
Package: msmtp
Version: 1.6.6-1
（略）
Depends: libc6 (>= 2.22), libgnutls30 (>= 3.5.6), libgsasl7 (>= 1.1), debconf (>= 0.5) | debconf-2.0, ucf
Recommends: ca-certificates
Suggests: msmtp-mta
Description-en: light SMTP client with support for server profiles
（以下略）
\end{verbatim}
 \normalsize
\end{frame}

\begin{frame}[containsverbatim]{mail-transport-agent (4/4)}
どのMTAがSMTPサーバも提供するかはパッケージのプログラム次第

 \begin{itemize}
  \item SMTPサーバも提供する -- exim4-daemon-light, postfix, nullmailer
	など
  \item SMTPサーバを提供しない（リレー専用） -- msmtp など
 \end{itemize}
\end{frame}

\section{caff と mail-transport-agent}
\begin{frame}[containsverbatim]{caff と mail-transport-agent (1/2)}
caff は署名済みの公開鍵をその人にメールで送信 

 \begin{itemize}
  \item 実行されたマシン上にある（ローカル）MTA を使ってメール送信している
  \item デフォルトの MTA だとローカル配送のみ
 \footnotesize
\begin{verbatim}
Subject: Mail delivery failed: returning message to sender
（略）
	 Mailing to remote domains not supported
\end{verbatim}
 \normalsize
	インターネットに出ていかない！ MTAの設定が必要
 \end{itemize}

\end{frame}

\begin{frame}{caff と mail-transport-agent (2/2)}
        もし自分の（uidに使っている）メールアドレスが自分の管理するメー
       ルサーバのドメインでないのなら

 \begin{itemize}
  \item メールサーバ（SMTPサーバ）なし
  \item リレー専用
 \end{itemize}

       なMTAを選べば設定がかんたん\\
       -- 例: msmtp-mta パッケージ

\pause
 \url{https://wiki.debian.org/msmtp}

 \begin{itemize}
  \item  個人の設定ファイル \textasciitilde/.msmtprc にGmailアカウントとSMTPサーバの情報を書いてリレー
 する方法が書かれている
  \item  パスワードをGPG鍵で暗号化して保存する方法も書かれている
 \end{itemize}
\end{frame}

\section{まとめ}
\begin{frame}{まとめ}
\begin{itemize}
 \item caff はローカルのMTAを使ってメール送信する
 \item ローカルのMTAはmail-transport-agent仮想パッケージを提供するパッケー
       ジをインストールすればよい
 \item マシンにはmail-transport-agent仮想パッケージを提供するパッケージのうち
       1つだけをインストールできる(Conflicts/Replaces/Provides)
 \item MTAのパッケージが2つに分かれていることもある（msmtp/msmtp-mta,
       qmail/qmail-runなど）
 \item MTA にはSMTPだけでなくsendmail互換インターフェースもある -- caff作業マシンが
       インターネット公開のメールサーバでないならSMTPは不要
 \item 自分が管理していないドメインのメールアドレスなら、その外部SMTPサー
       バにリレーするようにローカルのMTAをインストール（msmtp-mta パッケー
       ジがおすすめ）し設定するとよい
\end{itemize}
\end{frame}

\end{document}

;;; Local Variables: ***
;;; outline-regexp: "\\([ 	]*\\\\\\(documentstyle\\|documentclass\\|emtext\\|section\\|begin{frame}\\)\\*?[ 	]*[[{]\\|[]+\\)" ***
;;; End: ***
